\chapter*{Abstract}
\addcontentsline{toc}{chapter}{Abstract}

Neutrinoless double beta decay ($0\nu\beta\beta$) is a hypothetical nuclear process in which two neutrons turn into two protons and two electrons without the emission of any neutrino. Its observation would prove that lepton number is not conserved and that neutrinos are Majorana particles, i.e.\ their own antiparticles, providing direct evidence of physics beyond the Standard Model, with implications for the origin of neutrino masses and the matter-antimatter asymmetry of the Universe. LEGEND is an experimental programme searching for this decay in~${}^{76}$Ge with high-purity germanium (HPGe) detectors enriched in this isotope, which act at the same time as source and detector.

Since the decay, if it exists, is extremely rare, the search relies on pulse shape discrimination (PSD) to reject multi-site background events (MSE) while preserving single-site, signal-like events (SSE). The standard technique condenses each current pulse into a single parameter, the ratio A/E between the maximum current amplitude and the energy. This thesis investigates whether deep learning models exploiting the full time evolution of the pulse can match or improve on it.

The study is based on calibration data from two of the detectors used in LEGEND-200, in an energy window containing a peak dominated by SSE and one dominated by MSE. Two complementary approaches were implemented in PyTorch and trained on a GPU. The first is a supervised classifier built on the latent space of a Variational Autoencoder (VAE), trained on labels given by energy windows around the two peaks, whose purity is limited by the underlying Compton continuum. The second is a fully unsupervised Deep Temporal Clustering (DTC) model, which combines a convolutional-recurrent temporal autoencoder with a clustering layer and was tested with Euclidean, Complexity Invariant (CID) and Autocorrelation-based (ACF) distances.

The deep learning cuts are compared with A/E on the calibration data, through net survival fractions extracted from peak fits, the corresponding ROC curves, and the study of the events on which the methods agree or disagree. Finally, all the cuts are applied to $2\nu\beta\beta$-dominated physics data, never seen during training, to test their behaviour on a different type of data that spread over a wider energy range and that have precisely the same pulse shape of the $0\nu\beta\beta$ events that LEGEND is looking for.